
\chapter{Methodology}
\label{chap:methodology}

This chapter presents the algorithms that underpin our experimental study.
We begin by formalising goal-conditioned reinforcement learning in a learned
latent space (\S\ref{sec:problem}), followed by a description of the
JEPA-style world model used to construct this representation
(\S\ref{sec:lewm-bg}). 

Building on this foundation, we introduce a sequence of reinforcement
learning methods, each motivated by a concrete limitation of the previous
approach. Rather than presenting these methods in isolation, we develop them
incrementally, highlighting how each design decision addresses specific
failures encountered in practice. The methods are therefore organised both
chronologically and conceptually: curriculum SAC
(\S\ref{sec:curr-sac}), HER-enhanced SAC
(\S\ref{sec:her-sac}), hierarchical SAC with behavioural regularisation
(\S\ref{sec:joint}), HIQL adapted to latent world models
(\S\ref{sec:hiql-lewm}), and finally a flow-matching policy for
high-dimensional action generation (\S\ref{sec:hiql-flow}).

\medskip

The progression of methods can be understood as a sequence of increasingly
structured solutions to long-horizon, high-dimensional control in offline
settings.

\textbf{(i)~Latent World Model (LeWM).}
We first decouple representation learning from control by training a
JEPA-style world model (\S\ref{sec:lewm-bg}). The model consists of an
encoder and predictor trained with a spectral isotropy regulariser,
producing a latent space in which Euclidean distances are approximately
Gaussian and isotropic. This geometric structure is crucial: it enables
simple distance-based rewards and stabilises downstream learning.
Once trained, the encoder and predictor are frozen, and all subsequent
methods operate entirely within this latent space.

\textbf{(ii)~Curriculum SAC.}
Using the learned latent dynamics as a differentiable simulator, we next
train a flat goal-conditioned policy with Soft Actor-Critic
(\S\ref{sec:curr-sac}). To mitigate the difficulty of long-horizon tasks,
we introduce a curriculum over goal distances, gradually increasing task
complexity during training. Despite careful tuning—including action
penalties, reward shaping, and stable initialisation—this approach
struggles to scale. Without hindsight relabelling, the agent receives
little to no informative signal when goals are distant, resulting in slow
or unstable learning.

\textbf{(iii)~HER + SAC.}
To address this sparsity, we incorporate Hindsight Experience Replay
(\S\ref{sec:her-sac}). By relabelling trajectories with goals that were
actually achieved, every transition becomes a potential learning signal.
We further introduce an episodic replay buffer, explicit sparse and dense
reward formulations, and multiple goal-sampling strategies (curriculum,
fixed-gap, and distance-based). This significantly improves data
efficiency and broadens the effective training distribution. However,
the policy remains flat, and performance degrades on tasks requiring
multi-step planning.

\textbf{(iv)~Hierarchical SAC with Behavioural Regularisation.}
To better handle long-horizon reasoning, we introduce a hierarchical
structure (\S\ref{sec:joint}). A high-level policy proposes intermediate
waypoints, while a low-level SAC agent learns to reach them. The
high-level policy is trained using distance-based
advantage-weighted regression, and its behaviour is regularised using a
behaviour-cloning loss derived from an offline expert cache. This
stabilises early training and encourages realistic subgoal selection.
However, the approach remains sensitive to compounding model errors:
rollouts in the learned latent dynamics can drift over time, and the
action-conditioned $Q$-function exhibits high variance, limiting
reliability.

\textbf{(v)~HIQL + LeWM.}
To overcome the instability of online value learning in latent
imagination, we reformulate the problem as offline reinforcement learning
using Implicit Q-Learning (\S\ref{sec:hiql-lewm}). Instead of learning an
action-conditioned value function $Q(z, g, a)$, we learn an action-free
value function $V(z, \varphi(z,g))$ using expectile regression. A learned
low-dimensional goal embedding $\varphi$ compresses the high-dimensional
latent space, improving generalisation and stability. Crucially, this
formulation avoids querying out-of-distribution actions, making it well
suited to training on static datasets. World-model rollouts are only
introduced after an initial warm-up phase to prevent early-stage
divergence.

\textbf{(vi)~Flow-Matching Policy.}
Finally, we address a key limitation of standard actor-critic methods:
the inability to represent complex, multimodal action distributions. We
replace the unimodal Gaussian low-level policy with a conditional
flow-matching model (\S\ref{sec:hiql-flow}), which learns a velocity field
over actions and is trained via one-step distillation. This allows the
policy to represent diverse manipulation strategies and produce coherent
action sequences in high-dimensional spaces, without requiring explicit
likelihood computation.

\medskip

Taken together, these components form a unified framework that combines
latent world modelling, hierarchical abstraction, offline value learning,
and expressive policy representations. The resulting system is designed to
address the core challenges of long-horizon, high-dimensional,
goal-conditioned control in offline settings.

% ---------------------------------------------------------------------
\section{Problem Formulation}
\label{sec:problem}

We consider goal-reaching as a goal-conditioned Markov decision process
(GCMDP) $\mathcal{M}_g = (\mathcal{S}, \mathcal{A}, \mathcal{G}, p, r, \gamma)$,
where $\mathcal{S}$ denotes the observation space, $\mathcal{A} = \mathbb{R}^{d_a}$
is a continuous action space, $\mathcal{G} \subseteq \mathcal{S}$ is the set of
valid goals, $p(s' \mid s,a)$ defines the (unknown) environment dynamics,
$r(s,a,g)$ is a goal-conditioned reward function, and $\gamma \in [0,1)$ is the
discount factor. The objective is to learn a policy $\pi(a \mid s, g)$ that
maximises expected return for arbitrary goals $g$.

In our setting, observations are high-dimensional RGB images
($224 \times 224$), making direct learning in $\mathcal{S}$ both
computationally inefficient and statistically challenging. To address this,
we introduce a learned latent representation $\mathcal{Z} = \mathbb{R}^{d_z}$
with $d_z = 192$, obtained via an encoder $E_\xi : \mathcal{S} \to \mathcal{Z}$.
This induces a latent goal-conditioned MDP:
\begin{equation}
\widehat{\mathcal{M}}_g
= (\mathcal{Z}, \mathcal{A}, \mathcal{Z}, \hat{p}, \hat{r}, \gamma),
\end{equation}
where both states and goals are represented in the same latent space.

The latent dynamics are modelled by an action-conditioned predictor
\begin{equation}
\hat{p}(z' \mid z, a) \approx f_\psi(z, a),
\end{equation}
which deterministically predicts the next latent state. Together, the encoder
and predictor $(E_\xi, f_\psi)$ constitute the world model described in
\S\ref{sec:lewm-bg}. Once trained, this model is frozen and used as a
differentiable simulator for all downstream reinforcement learning methods.

Operating in latent space provides two key advantages. First, it removes the
need to model high-dimensional pixel observations, focusing instead on
task-relevant features. Second, as discussed in \S\ref{sec:lewm-bg}, the latent
space is structured such that Euclidean distances are meaningful, enabling the
use of simple distance-based reward functions for goal-reaching.

\medskip

\noindent
We use the following notation throughout the remainder of this chapter:
\medskip

\noindent
\begin{tabular}{ll}
$z, z' \in \mathcal{Z}$ & latent states \\
$g \in \mathcal{Z}$ & goal embedding (latent goal) \\
$a \in \mathcal{A}$ & action \\
$\pi(a \mid z, g)$ & goal-conditioned policy \\
$Q_\theta(z, g, a)$ & action-value function (critic) \\
$V_\phi(z, g)$ & state-value function \\
$f_\psi(z, a)$ & latent dynamics predictor \\
$E_\xi(s) \to z$ & encoder mapping observations to latent space (frozen) \\
\end{tabular}

\section{The Learned World Model}
\label{sec:lewm-bg}

This work builds upon an existing latent world model, \emph{LeWorldModel}
(LeWM), which we adopt as a fixed representation and dynamics module
throughout our experiments. We do not contribute to the design or training
of this model; rather, we use an open-source implementation trained on
expert trajectories from \textsc{OGBench}. This section briefly summarises
the key components of LeWM and highlights the properties that are critical
for downstream control.

LeWM follows a Joint-Embedding Predictive Architecture (JEPA)
paradigm, in which the model is trained to predict future latent
representations rather than reconstruct observations at the pixel level.
Given a sequence of observations, the model learns a compact latent space
and a predictor that captures environment dynamics directly in this space.

\paragraph{Architecture.}
The model consists of two components: an encoder $E_\xi$ and a predictor
$f_\psi$. The encoder is a Vision Transformer (ViT-Tiny) that maps each
input frame $s_t$ to a latent representation
\[
z_t = E_\xi(s_t) \in \mathbb{R}^{d_z}, \quad d_z = 192.
\]
Actions are embedded into the same latent dimension and provided as
conditioning inputs to the predictor.

The dynamics model $f_\psi$ is implemented as a multi-layer Transformer,
which takes a short history of latent states and actions and predicts the
next latent state. Conditioning on actions is achieved via adaptive
normalisation (AdaLN), allowing the model to modulate its predictions based
on control inputs.

\paragraph{Predictive objective.}
The model is trained to minimise a mean-squared prediction error in latent
space. Given a context window of past states and actions, the predictor
outputs a one-step prediction:
\begin{equation}
\mathcal{L}_{\mathrm{pred}}
= \mathbb{E}_{\tau \sim \mathcal{D}} \left[
\big\| f_\psi(z_{t-c:t-1}, a_{t-c:t-1}) - z_t \big\|_2^2
\right],
\label{eq:pred-loss-clean}
\end{equation}
where $z_t = E_\xi(s_t)$. This objective encourages the latent space to
capture the minimal information required for predicting future dynamics.

\paragraph{Preventing representation collapse.}
A key challenge in JEPA-style models is representation collapse, where all
inputs are mapped to a constant vector. LeWM addresses this using a
\emph{spectral isotropy regulariser} (SIGReg), which encourages the latent
distribution to match an isotropic Gaussian.

Rather than directly constraining the covariance, SIGReg matches the
empirical characteristic function of the latent distribution to that of a
standard normal along random projection directions. This promotes both
diversity and isotropy in the learned representation, preventing degenerate
solutions while avoiding the need for architectural tricks such as
stop-gradients or target networks.

The overall training objective is:
\begin{equation}
\mathcal{L}_{\mathrm{LeWM}}
= \mathcal{L}_{\mathrm{pred}} + \lambda_{\mathrm{sig}} \,
\mathcal{L}_{\mathrm{SIGReg}}.
\label{eq:lewm-total-clean}
\end{equation}

\paragraph{Usage in this work.}
After training, both the encoder $E_\xi$ and predictor $f_\psi$ are frozen
and used as a differentiable simulator in latent space. All reinforcement
learning methods described in the remainder of this chapter operate solely
on latent states $z \in \mathcal{Z}$ and predicted transitions
$\hat{z}_{t+1} = f_\psi(z_t, a_t)$.

\paragraph{Implications for control.}
Two properties of LeWM are particularly important for our setting.

\emph{First}, the predictor $f_\psi$ enables fast and differentiable
rollouts in latent space, allowing policies to be trained without
expensive interaction with the underlying simulator.

\emph{Second}, the combination of latent prediction and isotropy
regularisation induces a structured geometry in $\mathcal{Z}$. In
particular, Euclidean distances $\|z - g\|_2$ provide a meaningful and
approximately isotropic measure of similarity between states. This
property is central to our formulation of goal-conditioned rewards, where
distance in latent space serves as a dense proxy for task progress.

In all subsequent methods, we define rewards of the form
$r = -\|z_{\text{next}} - z_{\text{goal}}\|_2$,
which is justified by the geometric properties of the latent space.

% ---------------------------------------------------------------------
\section{Curriculum Soft Actor-Critic (SAC) in Latent Space}
\label{sec:curr-sac}

We begin with the simplest learning setup that leverages the learned
world model: a flat, goal-conditioned Soft Actor-Critic (SAC) agent
trained entirely in latent space. The goal of this stage is not to
achieve state-of-the-art performance, but to establish a baseline and
identify the core challenges of learning from self-generated data in a
long-horizon setting.

The agent interacts only with the latent dynamics model $f_\psi$,
treating it as a differentiable simulator. No expert data or hindsight
relabeling is used at this stage. Instead, the policy is trained from
scratch using trajectories generated by rolling out the current policy,
with a curriculum over goal difficulty to stabilise learning.

\paragraph{Policy and value networks.}
Both the actor and critic are parameterised as multilayer perceptrons
with two hidden layers of width $256$ and layer normalisation. The actor
outputs a Gaussian policy with diagonal covariance, followed by a
$\tanh$ squashing function to enforce bounded actions:
\begin{equation}
a = s_{\max}\,\tanh\!\big(\mu_\phi(z,g) + \sigma_\phi(z,g)\odot\epsilon\big),
\qquad \epsilon \sim \mathcal{N}(0,I),
\end{equation}
where $s_{\max}$ controls the action scale. The corresponding log-probability
accounts for the squashing transformation:
\begin{equation}
\log\pi_\phi(a \mid z,g)
= \log\mathcal{N}(\tilde a;\,\mu_\phi,\sigma_\phi^2)
- \sum_{i=1}^{d_a} \log\!\big(s_{\max}(1-\tanh^2 \tilde a_i) + \varepsilon\big).
\end{equation}

Careful initialisation is required for stability. In particular, output
layers are initialised with small weights and the initial policy
variance is kept low, ensuring that early rollouts remain within the
support of the learned dynamics model. The critic consists of twin
Q-functions $\{Q_{\theta_1}, Q_{\theta_2}\}$, following standard SAC
practice to reduce overestimation bias.

\paragraph{Latent environment and reward.}
The environment is defined directly in latent space using the learned
predictor. Given a current state $z_t$ and action $a_t$, the next state
is obtained as $\hat{z}_{t+1} = f_\psi(z_t, a_t)$. 

We define a dense, goal-conditioned reward based on Euclidean distance:
\begin{equation}
r_t = -\|z_{t+1} - g\|_2 - \lambda_a \|a_t\|_2,
\end{equation}
where the first term encourages progress toward the goal and the second
term regularises the magnitude of actions. Episodes terminate when the
agent reaches a neighbourhood of the goal:
\begin{equation}
d_t = \mathbb{1}\!\big[\|z_t - g\|_2 < \delta\big].
\end{equation}
The threshold $\delta$ is chosen based on the typical scale of latent
prediction error, ensuring that termination corresponds to meaningful
goal achievement.

\paragraph{Optimisation.}
Training follows the standard SAC objective. The critic is trained using
a soft Bellman backup:
\begin{equation}
\widehat{Q}_t = r_t + \gamma (1 - d_t)
\Big(\min_{i=1,2} Q_{\bar{\theta}_i}(z', g, a') - \alpha \log \pi_\phi(a' \mid z', g)\Big),
\end{equation}
and minimises the squared error to this target. The actor is updated by
maximising expected Q-values while encouraging entropy:
\begin{equation}
\mathcal{L}_\pi = \mathbb{E}\!\left[\alpha \log \pi_\phi(a \mid z,g)
- \min_{i=1,2} Q_{\theta_i}(z,g,a)\right].
\end{equation}
The temperature parameter $\alpha$ is learned to match a target entropy,
balancing exploration and exploitation.

\paragraph{Curriculum over goal difficulty.}
A key challenge in goal-conditioned learning is that distant goals
provide little useful signal early in training. To address this, we
introduce a curriculum over goal distances. At stage $k$, goals are
sampled from states that are $G_k$ steps ahead in expert trajectories,
and rollouts are limited to a matching horizon $T_{\max,k}$:
\begin{equation}
G_k \in \{1,2,4,8,16,24,32,40\}, \qquad T_{\max,k} = G_k.
\end{equation}
The agent progresses to the next stage once its success rate exceeds a
threshold, gradually increasing task difficulty.

\paragraph{Limitations.}
This approach relies entirely on on-policy data generated by the agent
itself. While the dense latent reward provides some guidance, it does not
fully resolve the challenge of long-horizon credit assignment. As the
goal distance increases, the agent rarely encounters successful
trajectories, and learning becomes unstable. In practice, performance
plateaus after a small number of curriculum stages.

This limitation highlights a fundamental issue: without additional data
reuse, the agent cannot efficiently learn from its failures. In the next
section, we address this by introducing hindsight goal relabeling, which
transforms every trajectory into a useful learning signal.

% ---------------------------------------------------------------------
\section{HER+SAC: Hindsight Relabelling and Reward Design}
\label{sec:her-sac}

The failure mode identified in \S\ref{sec:curr-sac} arises from a lack of
informative learning signal: the agent rarely reaches distant goals, and
therefore receives little feedback on how to improve. This limitation is
not specific to SAC, but reflects a broader issue in goal-conditioned
reinforcement learning—namely, that sparse success signals make
long-horizon learning intractable without additional data reuse.

To address this, we extend the baseline with Hindsight Experience Replay
(HER), transforming failed trajectories into useful supervision. In
addition, we introduce an episodic replay buffer and explicitly study two
reward formulations (sparse and dense), along with different strategies
for sampling training tasks.

\paragraph{Episodic replay buffer.}
Unlike the flat replay buffer used previously, we store complete
trajectories:
\begin{equation*}
\mathcal{B} = \big\{(z_{0:T_{\mathrm{ep}}},\,a_{0:T_{\mathrm{ep}}-1},\,z'_{0:T_{\mathrm{ep}}-1},\,g_{\mathrm{orig}})\big\}.
\end{equation*}
This structure is essential for HER, as it allows efficient sampling of
future states from within the same episode. In practice, episodes are
stored in pre-allocated GPU tensors for fast indexing and batch
construction.

\paragraph{Hindsight goal relabelling.}
For each sampled transition $(z_t, a_t, z'_t, g_{\mathrm{orig}})$, we
construct a relabelled goal by selecting a future state from the same
trajectory:
\begin{equation}
g_{\mathrm{rel}} =
\begin{cases}
z_{t+\Delta}, \quad \Delta \sim \mathrm{Unif}\{1, \dots, T_{\mathrm{ep}} - t - 1\}, & \text{with probability } p_{\mathrm{HER}}, \\
g_{\mathrm{orig}}, & \text{otherwise}.
\end{cases}
\label{eq:her-relabel-clean}
\end{equation}
Intuitively, this treats every visited state as a potential goal,
allowing the agent to learn from what it \emph{did} achieve, rather than
only what it intended to achieve. In our experiments, we use
$p_{\mathrm{HER}} = 0.8$.

\paragraph{Reward formulations.}
Given a relabelled goal $g_{\mathrm{rel}}$, rewards are recomputed
on-the-fly. We consider two distinct reward settings, evaluated in
separate experiments.

\textbf{Sparse reward.} The agent receives a reward only upon reaching
the goal:
\begin{equation}
r_t =
\mathbb{1}\!\big[\|z_{t+1} - g_{\mathrm{rel}}\|_2 < \delta\big] - 1
- \lambda_a \|a_t\|_2.
\label{eq:her-sparse-clean}
\end{equation}
This corresponds to a standard goal-reaching objective, where the agent
is penalised until success. The threshold $\delta$ matches the latent
termination condition defined earlier.

\textbf{Dense reward.} In a separate set of experiments, we instead use a
potential-based shaping reward:
\begin{equation}
r_t =
\mathrm{clip}\!\left(
\frac{\|z_t - g_{\mathrm{rel}}\|_2 - \|z_{t+1} - g_{\mathrm{rel}}\|_2}{\kappa},
-2,\,2
\right)
- \lambda_a \|a_t\|_2.
\label{eq:her-dense-clean}
\end{equation}
This reward measures progress toward the goal and can be interpreted as a
clipped difference of a potential function
$\Phi(z,g) = -\|z - g\|_2$. The scaling factor $\kappa$ controls the
magnitude of rewards and prevents individual transitions from producing
outliers.

These two reward formulations serve different purposes: the sparse
reward tests whether HER alone is sufficient to overcome credit
assignment, while the dense reward evaluates whether additional shaping
further improves learning efficiency.

\paragraph{Task sampling strategies.}
In addition to relabelling goals within trajectories, we vary how initial
state–goal pairs $(z_0, g)$ are constructed. We consider three sampling
strategies:

\begin{align}
\text{(curriculum)}\quad & (z_0, g) = (z_t,\, z_{t + G_k}), \quad G_k \uparrow,
\\
\text{(fixed-gap)}\quad & (z_0, g) = (z_t,\, z_{t + G}),
\\
\text{(pure-distance)}\quad & (z_0, g) = (z_t,\, z_{T_{\mathrm{ep}}-1}).
\end{align}

The curriculum strategy mirrors \S\ref{sec:curr-sac}, gradually
increasing goal difficulty. The fixed-gap setting removes this
adaptivity, while the pure-distance setting represents the most
challenging case, requiring the agent to reach the final state of a
trajectory. In this regime, HER implicitly provides intermediate
subgoals through relabelling, effectively inducing a curriculum from the
data.

\paragraph{Optimisation.}
The underlying learning algorithm remains unchanged: we optimise the
standard SAC objectives (Eqs.~\eqref{eq:sac-q}--\eqref{eq:sac-alpha}),
with rewards defined by either Eq.~\eqref{eq:her-sparse-clean} or
Eq.~\eqref{eq:her-dense-clean}. 

In practice, the dense reward leads to lower-variance updates and thus
requires fewer gradient steps per iteration, while the sparse reward
relies more heavily on HER to provide sufficient training signal.

\paragraph{Discussion.}
HER fundamentally changes the learning paradigm from purely on-policy
trial-and-error to off-policy data reuse. By extracting supervision from
every trajectory, it significantly improves sample efficiency and enables
learning on tasks that were previously infeasible.

However, while HER alleviates reward sparsity, it does not address the
core challenge of long-horizon reasoning. The policy remains flat, and as
task complexity increases, performance is still limited by the inability
to decompose problems into intermediate steps. This motivates the
introduction of hierarchical structure in the following section.

% ---------------------------------------------------------------------
\section{Hierarchical SAC with Expert-Guided Stabilisation}
\label{sec:joint}

While hindsight relabelling significantly improves sample efficiency, it
does not resolve a more fundamental limitation: the policy remains
\emph{flat}. As task horizons increase, the agent must reason over long
sequences of actions, and the action-conditioned value function
$Q(z,g,a)$ becomes increasingly difficult to estimate accurately from
world-model rollouts alone. In particular, distant goals induce high
variance in both value estimates and policy updates, leading to unstable
training.

To address this, we introduce a hierarchical structure that decomposes
long-horizon tasks into shorter, more tractable subproblems. In addition,
we incorporate offline expert data as a stabilising signal during
training.

\paragraph{Hierarchical decomposition.}
We factor the policy into two levels:
\begin{equation}
\pi^{\mathrm{HL}}(z_{\mathrm{sub}} \mid z, g_{\mathrm{ult}}),
\qquad
\pi^{\mathrm{LL}}(a \mid z, z_{\mathrm{sub}}),
\end{equation}
where $g_{\mathrm{ult}}$ is the ultimate goal and
$z_{\mathrm{sub}} \in \mathcal{Z}$ is a latent subgoal predicted by the
high-level policy. Intuitively, the high-level actor proposes an
intermediate waypoint, and the low-level actor is trained to reach this
subgoal using primitive actions.

The high-level policy operates at a lower temporal resolution, predicting
subgoals approximately $G$ steps into the future, while the low-level
policy executes at every timestep.

\paragraph{Low-level policy learning.}
The low-level policy is trained using SAC, as in
\S\ref{sec:curr-sac}, but with the subgoal $z_{\mathrm{sub}}$ replacing
the global goal. To stabilise training, we augment the actor objective
with a behaviour cloning (BC) term derived from an offline expert
dataset $\mathcal{D}_{\mathrm{exp}}$:
\begin{equation}
\mathcal{L}_\pi^{\mathrm{LL}} =
\mathbb{E}\!\left[\alpha \log \pi^{\mathrm{LL}}(a \mid z, z_{\mathrm{sub}})
- \min_{i=1,2} Q_{\theta_i}(z, z_{\mathrm{sub}}, a)\right]
+ \beta_{\mathrm{bc}}\,
\mathbb{E}_{(z,g,a^*) \sim \mathcal{D}_{\mathrm{exp}}}
\!\left[\|\pi^{\mathrm{LL}}(z,g) - a^*\|_2^2\right].
\end{equation}
This regularisation anchors the policy to the support of expert actions,
reducing the likelihood of generating out-of-distribution actions that
would lead to compounding errors in the world model.

\paragraph{Reward design.}
Rewards are defined as in \S\ref{sec:her-sac}, but with the goal replaced
by the current subgoal $z_{\mathrm{sub}}$. In the dense reward setting,
we disable hindsight relabelling, as trivial relabelling (e.g.,
$g_{\mathrm{rel}} = z_{t+1}$) would artificially maximise the shaping
reward and introduce degenerate learning signals.

\paragraph{High-level policy learning.}
The high-level policy is trained using advantage-weighted regression
(AWR), with a distance-based notion of advantage. For transitions sampled
from trajectories, we define:
\begin{equation}
A_t^{\mathrm{HL}} =
\|z_t - g_{\mathrm{ult}}\|_2 - \|z_{t+G} - g_{\mathrm{ult}}\|_2,
\end{equation}
which measures the improvement in goal proximity over $G$ steps. This is
converted into a weight:
\begin{equation}
w_t = \min\!\big(\exp(\beta_{\mathrm{HL}} A_t^{\mathrm{HL}}),\, w_{\max}\big),
\end{equation}
and the high-level objective becomes:
\begin{equation}
\mathcal{L}_{\mathrm{HL}} =
\mathbb{E}\!\left[
w_t \,\big\|\pi^{\mathrm{HL}}(z_t, g_{\mathrm{ult}}) - z_{t+G}\big\|_2^2
\right].
\end{equation}
This encourages the high-level policy to predict subgoals that are both
reachable and make meaningful progress toward the ultimate goal.

\paragraph{Expert data integration.}
To further improve stability, we mix transitions from the expert dataset
$\mathcal{D}_{\mathrm{exp}}$ with those generated by the agent. Each
training batch contains a fraction of expert transitions, controlled by
a mixing parameter. This provides a strong initial training signal and
helps prevent early collapse of the policy.

At one extreme, using only expert data reduces to pure imitation
learning; at the other, relying solely on self-generated data recovers
the HER+SAC setting. In practice, an intermediate mixture yields the best
performance.

\paragraph{Optional latent projection.}
In some experiments, we project latent states to a lower-dimensional
space before feeding them to the reinforcement learning components. This
reduces variance in value estimation and sharpens the distance-based
signals used by the high-level policy, while leaving the world model
itself unchanged.

\paragraph{Limitations.}
This hierarchical formulation substantially improves performance on
longer-horizon tasks by decomposing them into manageable subproblems.
However, training remains fundamentally on-policy and reliant on
world-model rollouts. As a result, it is still sensitive to model errors
and requires careful balancing between exploration, imitation, and
stability.

These challenges motivate the transition to a fully offline formulation,
introduced in the next section, which avoids direct reliance on
action-conditioned value estimation altogether.

% ---------------------------------------------------------------------
\section{HIQL in the LeWM Latent Space}
\label{sec:hiql-lewm}

The hierarchical SAC formulation in \S\ref{sec:joint} alleviates
long-horizon planning by introducing temporal abstraction, but it still
relies on an action-conditioned critic $Q(z,g,a)$. In practice, this
remains a major source of instability: estimating $Q$ requires reasoning
over both state and action spaces, and small errors in the world model
compound when evaluating out-of-distribution actions.

To address this, we move to an \emph{offline} formulation based on
Implicit Q-Learning (IQL), and adopt its hierarchical extension (HIQL)
in the LeWM latent space. The key idea is to \emph{eliminate explicit
dependence on actions in the value function}, replacing $Q(z,g,a)$ with
an action-free value function $V(z, \varphi(z,g))$. This removes the need
to evaluate unseen actions and significantly improves stability.

\paragraph{Model components.}
The world model $(E_\xi, f_\psi)$ remains fixed. The trainable
components are:
\begin{equation}
\varphi_\eta(z,g) \in \mathbb{R}^{d_r}, \qquad d_r = 10,
\qquad \|\varphi_\eta(z,g)\|_2 = \sqrt{d_r},
\end{equation}
\begin{equation}
V_{\phi_1}, V_{\phi_2} : \mathcal{Z} \times \mathbb{R}^{d_r} \to \mathbb{R},
\qquad
\pi^{\mathrm{HL}}(\varphi \mid z,g), \quad
\pi^{\mathrm{LL}}(a \mid z,\varphi).
\end{equation}

Here, $\varphi_\eta(z,g)$ is a learned low-dimensional representation of
the goal, which compresses the high-dimensional latent space into a more
manageable form. This is particularly important because the LeWM latent
space ($d_z=192$) is significantly overparameterised for value learning.
Projecting onto a $10$-dimensional sphere enforces a structured and
bounded representation, improving generalisation and stability.

\paragraph{Value learning via expectile regression.}
Instead of learning a $Q$-function, HIQL learns a value function that
implicitly captures the maximum over actions. Given a subgoal
$z_{\mathrm{sub}} = z_{t+k}$ sampled from the dataset, we define a dense
reward:
\[
\hat r_t = -\|z_{t+1} - z_{\mathrm{sub}}\|_2,
\]
consistent with the latent geometry induced by LeWM.

We construct target values:
\begin{align}
q_i &= \hat r_t + \gamma\,V_{\bar\phi_i}\!\big(z_{t+1},\,\bar\varphi(z_{t+1}, z_{\mathrm{sub}})\big),
\end{align}
and define an advantage-like quantity:
\begin{align}
A_t = \frac{1}{2}\sum_i q_i - \frac{1}{2}\sum_i V_{\bar\phi_i}\!\big(z_t,\,\bar\varphi(z_t, z_{\mathrm{sub}})\big).
\end{align}

The value function is trained using \emph{expectile regression}:
\begin{align}
\ell_\tau(u) &= \big|\tau - \mathbb{1}[u < 0]\big|\,u^2, \qquad \tau = 0.7, \\
\mathcal{L}_V &= \mathbb{E}\!\left[\sum_{i=1}^{2}
\ell_\tau\!\Big(q_i - V_{\phi_i}\!\big(z_t, \varphi_\eta(z_t, z_{\mathrm{sub}})\big)\Big)\right].
\end{align}

Intuitively, expectile regression biases the value function toward the
\emph{upper tail} of the observed returns. This serves as a stable proxy
for the $\max_a Q(z,g,a)$ operator in classical Q-learning, without ever
evaluating actions explicitly. As a result, the method avoids
overestimation errors on out-of-distribution actions—a key failure mode
in offline RL.

\paragraph{Policy learning via advantage-weighted regression.}
Both the low-level and high-level policies are trained using
advantage-weighted regression (AWR), which encourages imitation of
actions and subgoals that lead to high-value outcomes.

For the low-level policy:
\begin{align}
A^{\mathrm{LL}}_t &=
\frac{1}{2}\sum_i \Big[
V_{\phi_i}\!\big(z_{t+1}, \varphi(z_{t+1}, z_{\mathrm{sub}})\big)
- V_{\phi_i}\!\big(z_t, \varphi(z_t, z_{\mathrm{sub}})\big)
\Big], \\
\mathcal{L}_{\pi^{\mathrm{LL}}} &=
-\mathbb{E}\!\left[
\min\big(\exp(\alpha_{\mathrm{LL}} A^{\mathrm{LL}}_t), 100\big)
\log \pi^{\mathrm{LL}}\!\big(a_t \mid z_t, \varphi(z_t, z_{\mathrm{sub}})\big)
\right].
\end{align}

For the high-level policy:
\begin{align}
A^{\mathrm{HL}}_t &=
\frac{1}{2}\sum_i \Big[
V_{\phi_i}\!\big(z_{t+k}, \varphi(z_{t+k}, g_{\mathrm{ult}})\big)
- V_{\phi_i}\!\big(z_t, \varphi(z_t, g_{\mathrm{ult}})\big)
\Big], \\
\mathcal{L}_{\pi^{\mathrm{HL}}} &=
-\mathbb{E}\!\left[
\min\big(\exp(\alpha_{\mathrm{HL}} A^{\mathrm{HL}}_t), 100\big)
\log \pi^{\mathrm{HL}}\!\big(\varphi(z_t, z_{\mathrm{sub}}) \mid z_t, g_{\mathrm{ult}}\big)
\right].
\end{align}

The key idea is that both policies learn by \emph{reweighting behaviour}
according to value improvement: transitions that lead to greater progress
toward the goal are exponentially upweighted.

\paragraph{Offline training with optional imagination.}
HIQL is fundamentally an offline method, trained on fixed datasets of
transitions. However, to leverage the learned world model, we introduce a
two-phase training procedure:
\begin{align}
\textbf{Phase 1 (offline warm-up):} &\quad \text{train only on real data}, \\
\textbf{Phase 2 (augmentation):} &\quad \text{mix real and synthetic data}.
\end{align}

Synthetic transitions are generated by rolling out the low-level policy
through the latent dynamics model $f_\psi$. Subgoals are still sampled
from real trajectories, ensuring that the imagined data remains grounded
in the dataset distribution.

\paragraph{Stability mechanisms.}
To stabilise training, we maintain target networks for both the value
function and goal representation:
\[
\bar\phi \leftarrow (1 - \tau)\bar\phi + \tau \phi, \qquad
\bar\eta \leftarrow (1 - \tau)\bar\eta + \tau \eta,
\]
with $\tau = 0.005$.

\paragraph{Discussion.}
By removing the action-conditioned critic and adopting an offline
learning paradigm, HIQL eliminates the primary source of instability
observed in earlier methods. The combination of expectile regression and
advantage-weighted regression enables stable learning from static data,
while still supporting hierarchical reasoning.

However, the low-level policy remains a unimodal Gaussian, which limits
its ability to represent complex, multimodal action distributions
commonly encountered in manipulation tasks. This limitation motivates the
final extension presented in the next section.

% ---------------------------------------------------------------------
\section{Flow-Matching Low-Level Policy}
\label{sec:hiql-flow}

The HIQL formulation in \S\ref{sec:hiql-lewm} removes instability arising
from action-conditioned value estimation, but the low-level policy
remains a unimodal Gaussian. This assumption is restrictive in
manipulation tasks, where multiple distinct actions may be equally valid
in the same state (e.g., approaching an object from different
directions). As a result, the Gaussian policy tends to average over
modes, producing suboptimal or physically implausible actions.

To address this limitation, we replace the Gaussian low-level actor with
a \emph{flow-matching policy} that can represent complex, multimodal
action distributions. The policy consists of two components: (i) a
conditional velocity field trained via behaviour cloning, and (ii) a
distilled one-step policy trained using advantage-weighted regression.
All other components (value function, high-level policy, goal
representation, and training procedure) remain unchanged.

\paragraph{Conditional flow matching.}
We model the action distribution implicitly as a continuous
transformation from Gaussian noise to expert actions. Specifically, we
define a velocity field $u_\theta$ that describes how a sample evolves
over time from an initial noise distribution to a target action.

Following the flow-matching formulation, we sample
\[
x_0 \sim \mathcal{N}(0,I), \qquad x_1 = a^*, \qquad t \sim \mathrm{Unif}[0,1],
\]
and construct an interpolation:
\[
x_t = (1 - t)\,x_0 + t\,x_1.
\]
The velocity field is trained to predict the direction from noise to the
target action:
\begin{equation}
\mathcal{L}_{\mathrm{BC}}(\theta)
= \mathbb{E}\!\left[
\big\| u_\theta(z, \varphi, x_t, t) - (x_1 - x_0) \big\|_2^2
\right].
\label{eq:flow-bc-clean}
\end{equation}

Intuitively, this learns a continuous vector field that transports a
simple Gaussian distribution into the complex, multimodal distribution
of expert actions, conditioned on the current state and goal.

\paragraph{Action generation via integration.}
To generate an action, we simulate this flow starting from Gaussian
noise. Using a simple Euler integration scheme with $S$ steps:
\begin{equation}
x_{i+1} = x_i + \frac{1}{S}\,u_\theta(z, \varphi, x_i, i/S),
\qquad x_0 = \epsilon \sim \mathcal{N}(0,I),
\end{equation}
and the final action is:
\begin{equation}
\hat{a}_{\mathrm{flow}} = \mathrm{clip}(x_S, -s_{\max}, s_{\max}).
\label{eq:flow-euler-clean}
\end{equation}

This procedure allows the policy to generate diverse actions depending
on the initial noise sample, naturally capturing multimodal behaviour.

\paragraph{Advantage-weighted distillation.}
While the flow model is expressive, iterative integration is
computationally expensive at inference time. To address this, we train a
second network $\mu_\psi$ to approximate the output of the flow in a
single forward pass.

This network takes a noise sample $\epsilon$ and predicts an action
directly. It is trained using advantage-weighted regression to imitate
the flow output:
\begin{equation}
\mathcal{L}_{\mathrm{dist}}(\psi)
= \mathbb{E}\!\left[
\min\big(\exp(\alpha_{\mathrm{LL}} A^{\mathrm{LL}}_t), 100\big)
\|\mu_\psi(z, \varphi, \epsilon) - \hat{a}_{\mathrm{flow}}\|_2^2
\right].
\label{eq:flow-dist-clean}
\end{equation}

The advantage weights ensure that the distilled policy focuses on
high-value behaviours, effectively selecting the most useful modes of
the action distribution.

\paragraph{Training objective.}
The two components are trained jointly:
\begin{equation}
\mathcal{L}^{\mathrm{LL}}_{\mathrm{flow}}
= \mathcal{L}_{\mathrm{BC}}(\theta)
+ \mathcal{L}_{\mathrm{dist}}(\psi).
\label{eq:flow-total-clean}
\end{equation}

The behaviour cloning loss uses all available expert transitions and
learns the full action distribution, while the distillation loss shapes
the final policy toward high-advantage regions.

\paragraph{Inference and practical considerations.}
At evaluation time, the agent uses only the distilled policy
$\mu_\psi$, requiring a single forward pass. This makes inference as
efficient as the Gaussian actor used in \S\ref{sec:hiql-lewm}, while
retaining the ability to model multimodal action distributions.

\paragraph{Discussion.}
This formulation resolves a key limitation of previous methods: the
inability of unimodal policies to represent diverse manipulation
strategies. By combining expressive generative modelling with
advantage-weighted distillation, the policy can capture multiple valid
behaviours while remaining efficient at deployment.

This completes the progression of methods in this chapter, culminating
in a stable, hierarchical, and expressive framework for offline
goal-conditioned control in learned latent spaces.

% ---------------------------------------------------------------------
\section{Implementation Notes}
\label{sec:impl}

\paragraph{Hardware.} All experiments are performed on a single A100
NVIDIA GPU. The replay buffers are GPU-resident; for the largest
configuration (\S\ref{sec:curr-sac}) the flat buffer occupies
approximately $1.5$\,GB.

\paragraph{Optimisers.} Adam with learning rate $3\!\times\!10^{-4}$
for all RL networks (actor, critic/value, temperature, goal-rep,
flow), and AdamW with learning rate $5\!\times\!10^{-5}$ for the
world model. Gradient clipping at $1.0$ is applied to every loss.

\paragraph{Target networks.} Polyak averaging with $\tau = 0.005$ is
applied to all critics, values and the HIQL goal-rep
$\varphi$. Target updates are performed once per gradient step.

\paragraph{Batching and parallelism.} All RL methods use a batch size
of $256$ and $40$ gradient updates per outer iteration (reduced to
$20$ in the dense reward setting of \S\S\ref{sec:her-sac}--\ref{sec:joint}).
World-model roll-outs are vectorised across $50$ parallel
environments.

\paragraph{Constants.} $\gamma = 0.99$ throughout. The done-threshold
$\delta_d = 2.0$ and the dense-reward scale $\kappa = 1.6$ are
calibrated from the median single-step predictor drift in the LeWM
latent space; the corresponding diagnostic is reported in
\S\ref{chap:experiments}. \todo{confirm $\delta_d$ and $\kappa$
match the values used in final runs.}


